% !TEX program = xelatex
% VLN-001 讲义笔记：Quantization by Branes and Geometric Langlands
% 由 13_lecture_notes/VLN-001_witten2021_branes_gl 的语料整理；规范见 04_playbook/24_视频讲义提取与笔记规范_v1.md
\documentclass[11pt,a4paper]{ctexart}
\usepackage[margin=2.35cm]{geometry}
\usepackage{amsmath,amssymb,amsthm}
\usepackage{booktabs}
\usepackage{enumitem}
\usepackage{xcolor}
\usepackage{array}
\usepackage[hidelinks]{hyperref}

\setlist{itemsep=0.15em,topsep=0.3em,leftmargin=1.5em}
\linespread{1.12}

% 讲义原文引用块：[t=hh:mm:ss]
\newcommand{\slq}[2]{%
  \begin{quote}\small\itshape #2\par
  \raggedleft\normalfont\footnotesize\textcolor{gray}{[t=#1]}\end{quote}}
% 整理者按
\long\def\zm#1{\par\smallskip\noindent{\color{black!60}整理者按：#1}\par\smallskip}
\newcommand{\Bcc}{\mathcal{B}_{\mathrm{cc}}}
\newcommand{\Loc}{\mathcal{M}_H}
\newcommand{\Bun}{\mathrm{Bun}_G}
\newcommand{\op}{\mathrm{op}}
\newcommand{\opb}{\overline{\mathrm{op}}}
\newcommand{\Hom}{\operatorname{Hom}}

\title{\bfseries 讲义笔记：Quantization by Branes and Geometric Langlands}
\author{Edward Witten（IAS）\\[2pt]
\small 整理：数译 · 视频讲义专栏 VLN-001}
\date{报告日期 2021-02-03 \quad·\quad 笔记版本 2026-11-27（B 级：未做系统抽核）}

\ctexset{section/format={\Large\bfseries\raggedright},subsection/format={\large\bfseries\raggedright}}
\begin{document}
\maketitle
\vspace{-1.2em}
\begin{center}\small\color{black!60}
Mini-Conference on the Geometric Langlands Correspondence · 02:04:53
\end{center}

\section*{0\quad 元数据摘要}
\addcontentsline{toc}{section}{0 元数据摘要}

\begin{center}\small
\begin{tabular}{@{}p{3.1cm}p{12.2cm}@{}}
\toprule
来源 & Bilibili \texttt{BV1b34y1Q7Ah}，分 P 4（该 P 直链 \texttt{?p=4}）\\
语言 & 讲义英文；音轨英文（Q\&A 夹杂提问者口音）\\
讲义规模 & 原 deck \textbf{228} 页；讲者实际翻到 \textbf{201} 页；逐条动画折叠后 \textbf{109} 个内容块\\
提取 & 1 fps 抽帧（915 帧）→ 页面分割 → 视觉转录（逐张读图，公式转 LaTeX）→ Whisper small.en 转写 → 按时间轴合并\\
权利 & 讲义原文与公式归讲者；本笔记的中文叙述、结构与时间轴为本仓自产（CC BY 4.0）；deck 页图像不入仓\\
\bottomrule
\end{tabular}
\end{center}

\section{导读}

这场报告的目的不是重新证明几何 Langlands，而是\textbf{给 EFK 的一套算子代数找物理原型}。
EFK（Etingof–Frenkel–Kazhdan，arXiv:1908.09677）在 $\mathcal{M}(G)$ 的 $\mathrm{L}^2$ 半密度上造了一族算子：
Hitchin 哈密顿量的量子化版本、以及 Hecke 算子；他们证明了若干对偶定理，并留下若干猜想。
问题在于：这些算子\textbf{为什么}是这些算子，谱\textbf{为什么}是那个谱。

Witten 给的回答是：把 $\Loc$ 看成实辛流形（$\omega=\operatorname{Im}\Omega$），用他和 Gukov 的
\textbf{量子化=边界条件}（brane quantization）方案，把 EFK 的 Hilbert 空间解释成 $A$-模型里的
$\Hom(\mathcal{B},\Bcc)$；于是"电-磁对偶"自动给出"谱的对应"——EFK 猜的
$\mathrm{L}^2(\mathcal{M}(G))$ 与 $\Hom(L_{\opb},L_{\op})$ 之间的等价，正是 $T$-对偶/镜像对称的一个实例。
最后他把这套机制推到 Hecke 算子的特征值：$SL_2$ 情形下，特征值就是实 oper 上的一对截面配对
$(\overline{j}(v),j(w))$，并且因为 $SL_2$ 的中心是 $\{\pm1\}$，还要带上一个 $\pm$。

一句话：\textbf{EFK 的猜想 = 一次量子化的谱对应，而量子化在 $A$-模型里就是"挑一个 brane"。}

\section{论证线（不读正文也能看全局）}

\begin{enumerate}[label=\arabic*.]
\item（00:01）几何 Langlands 的规范论版本：$\mathcal{N}=4$ 超对称规范论配合 $G\leftrightarrow G^\vee$ 的电-磁对偶，
耦合常数 $\tau$ 与 $\tau^\vee$ 成反演关系。
\item（00:04–00:15）边界条件构成范畴，线算子就是其上的函子；对偶把 $G$ 侧的 't Hooft 线算子（几何 Hecke）
对到 $G^\vee$ 侧的 Wilson 线算子，二者的合成代数一致，这就是几何 Satake。
\item（00:15–00:34）把上述一切翻译成镜像对称：$A$-模型在 $\Loc(G)$ 上 $\leftrightarrow$ $B$-模型在 $\Loc(G^\vee)$ 上；
$A$ 侧有一个特殊 brane $\Bcc$，它的自同态把 $A$-branes 变成 $\Bun$ 上的 $\mathcal{D}$-模；
$B$ 侧 $\Bcc$ 的对偶是 opers 的结构层（Beilinson–Drinfeld）。
\item（00:35–00:51）区分形变量子化与量子化（$S^2$ 的例子），然后给出 Gukov–Witten 判据：
$L$ 对 $\operatorname{Re}\Omega$ 是辛的 $\Rightarrow$ $\Hom(\mathcal{B},\Bcc)$ 是 $L$ 的量子化；
要成为 Hilbert 空间还需反全纯对合 $\tau$，$\langle\psi,\chi\rangle=(\Theta_\tau\psi,\chi)$。
\item（00:52–01:15）EFK 的 $\mathrm{L}^2(\mathcal{M}(G))$ 就是 $\Loc(G)\cong T^*\mathcal{M}(G)$ 的量子化；
可观测量是 $\mathcal{A}\otimes\overline{\mathcal{A}}$（Hitchin 哈密顿量及其共轭）；
对偶把它送到 $\Hom(L_{\opb},L_{\op})$，正定性 $\Leftrightarrow$ 交点是实 opers 且横截。
\item（01:16–01:26）Hecke 算子 $T_R(p)$ 与量子化的 Hitchin 哈密顿量交换（线算子可以滑动）。
\item（01:26–01:41）对偶预测 $T_R(p)$ 的特征值等于 $W_R(p)$ 在实 oper 上的值，即 $(\overline{j}(v),j(w))$；
"角"数据 $v,w$ 在对偶两侧必须相同，EFK 已用代数几何公式验证；$SL_2$ 需带 $\pm$。
\end{enumerate}

\section{出发点：Kapustin–Witten 的规范论表述}

几何 Langlands 在这里不是代数几何命题，而是一个四维场论的性质：取 $\mathcal{N}=4$ 超对称 Yang–Mills，
规范群 $G$ 或 Langlands–GNO 对偶群 $G^\vee$。

\slq{00:01:04}{The theory based on $G$ has a ``gauge coupling constant'' $e$ and topological angle $\theta$,
which combine to a complex parameter
\[
\tau=\frac{\theta}{2\pi}+\frac{4\pi}{e^{2}} .
\]
There is an elementary symmetry $\tau\cong\tau+1$. \dots\ The claim of ``electric-magnetic duality''
\dots\ is that the theories based on $G$ and on $G^{\vee}$ are equivalent under
\[
\tau^{\vee}=\frac{-1}{n_{\mathfrak{g}}\tau},
\]
$(n_{\mathfrak{g}}=1,2,\text{ or }3$ depending on $G$.)}

\zm{$\tau$ 的归一化与 $n_{\mathfrak{g}}$ 是 Witten 的记号习惯（$n_{\mathfrak{g}}$ 是 $G$ 的"层级"因子，
来自单连通性/中心），换文献时数值可能差一个因子。这句"$\tau\cong\tau+1$"下面还会用到：对偶与
$\tau\mapsto\tau+1$ 一起生成 $SL_2(\mathbb{Z})$ 作用。}

关键在于：这个对偶是\textbf{同一个场论的两种描述}，所以它作用在\textbf{一切}拓扑性质上——
包括边界条件与线算子。这正是后面所有内容的来源。

\section{边界条件、线算子与几何 Satake}

\subsection{场论给流形配什么}

\slq{00:04:50}{In general, quantum field theory in dimension $d$ associates a number to a $d$-manifold,
a vector space (``the space of physical states'') to a $d-1$-manifold, and a category
(``the category of boundary conditions'') to a $d-1$-manifold.}

报告只考虑 $M=\Sigma\times C$（$C$ 是固定的复曲线），于是固定的 $C$ 给出一个\textbf{范畴} $\mathcal{C}$，
其对象是 $\Sigma$ 的边界条件。两个边界条件 $\mathcal{B}_1,\mathcal{B}_2$ 之间：

\slq{00:06:49}{$\Hom(\mathcal{B}_1,\mathcal{B}_2)$ is the vector space that the quantum field theory
assigns to this picture. (Physicists usually don't use this terminology.)}

\zm{"这个图"指 $\Sigma$ 的一条边界上放 $\mathcal{B}_1$、另一条上放 $\mathcal{B}_2$、中间是场论路径积分。
物理学家一般说"在两个边界条件之间插入一个态"，数学家说 $\Hom$。这一步把图中的"线"变成代数里的"元素"。}

于是立刻拿到代数结构：

\slq{00:08:02}{With this definition, $\Hom(\mathcal{B},\mathcal{B})$ is an algebra for every $\mathcal{B}$,
and moreover $\Hom(\mathcal{B},\mathcal{B}')$ is a left $\Hom(\mathcal{B}',\mathcal{B}')$ module, etc.}

\subsection{线算子 = 函子}

\slq{00:10:00}{In a quantum field theory defined on a manifold $M$, a line operator is some sort of
modification of the definition of the theory along an embedded one-manifold $K\subset M$. The natural
line operators at $\Psi=\infty$ are what physicists call ``Wilson operators'' \dots\ and the natural
line operators at $\Psi=0$ are what physicists call ``'t Hooft operators.''}

\zm{$\Psi$ 是 Kapustin–Witten 的"扭转参数"。$\Psi=\infty$ 与 $\Psi=0$ 对应两种不同的拓扑扭转（$A$ 与 $B$），
所以同一个物理对象在两边的"自然"算子不同：Wilson 在 $\Psi=\infty$ 侧，'t Hooft 在 $\Psi=0$ 侧。}

报告里对线算子给出的精确说法是\textbf{函子}而不是算子：

\slq{00:10}{But for us, a line operator is a functor from the category of boundary conditions to itself.
We can understand that statement from this picture: [图 a) 一条线算子 $T$ 沿边界 $\mathcal{B}$ 放置，得到新边界条件 $T\mathcal{B}$；
b) 边界条件 $\mathcal{B},\mathcal{B}',\mathcal{B}''$ 串在一条 $T$ 上，说明 $T$ 如何作用于态射。]}

\zm{这张图是整场报告的"语法核心"：$T$ 沿边界走 = 作用在对象；$T$ 跨过若干边界条件 = 作用在态射。
后面"$T$ 与 $W$ 交换"、"$T$ 与 Hitchin 哈密顿量交换"，都是这张图的推论，不是独立计算。}

\subsection{对偶两侧：$T(p)$ 与 $W(p)$，几何 Satake}

\slq{00:15:35}{So the two dual categories $\mathcal{C}$ and $\mathcal{C}^{\vee}$ are equipped with dual
families $T(p)$ and $W(p)$ of functors, parametrized by the choice of $p\in C$ \dots
$T(p)$ corresponds to the usual Hecke functors of geometric Langlands, while $W(p)$ corresponds to
its dual \dots\ The duality says that the algebra of compositions $T_R(p)T_{R'}(p)$ of Hecke operators
at the same point is the same as the corresponding algebra of Wilson operators $W_R(p)W_{R'}(p)$,
which (on elementary grounds) is the tensor algebra of representations of $G^{\vee}$:
\[
W_R(p)W_{R'}(p)=W_{R\otimes R'}(p),
\]
where $W_{R\otimes R'}(p)$ can be expanded as a sum of irreducibles.
(This fact is known as the geometric Satake correspondence.)}

\zm{这段值得单独记住：\textbf{几何 Satake 在这里不是被证明的，而是"电-磁对偶 + Wilson 算子的张量代数"的直接后果}。
$W_R W_{R'}=W_{R\otimes R'}$ 在物理上"平凡"（把两个表示张起来做和乐），数学上却是 Satake 等价的核心。
Witten 自己用 "on elementary grounds" 一句带过，报告的重量放在下一步。}

\section{镜像对称的翻译与 canonical coisotropic brane}

\subsection{几何 Langlands = 镜像对称}

\slq{00:20}{Geometric Langlands duality can be understood for many purposes as a mirror symmetry
between the $A$-model of $\Loc(G)$ and the $B$-model of $\Loc(G^{\vee})$.}

\zm{$\Loc(G)$ 是 $G$-Higgs 丛的模空间，它同时是 $T^*\mathcal{M}(G)$ 的双有理模型（$\mathcal{M}(G)$ 是全纯 $G$-丛的模空间）。
"$A$-模型在 $\Loc(G)$ 上"意味着取复辛形式 $\Omega$、把 $\operatorname{Im}\Omega$ 当实辛形式；
"$\omega=\operatorname{Im}\Omega$"这个选取是后面所有正定性讨论的根据。}

\subsection{$\Bcc$ 与形变量子化}

\slq{00:29:01}{With this data one can construct a ``canonical coisotropic $A$-brane'' $\Bcc$, whose
support is all of $Y$. It has the property that, roughly, $\Hom(\Bcc,\Bcc)$ is related to deformation
quantization of the algebra of holomorphic functions on $Y$. More specifically, \dots\ if
$Y=T^*W$ \dots\ then $\Hom(\Bcc,\Bcc)$ is the sheaf of holomorphic differential operators on $W$,
acting on $K^{1/2}\to W$.}

\zm{"支撑是全空间但可以带 $\mathcal{B}$-场曲率"——这正是 coisotropic 与 Lagrangian 的区别。
$K^{1/2}$ 的出现是半密度，和 EFK 用 $\mathrm{L}^2$ 半密度是同一件事，后面会重合。}

落到几何 Langlands：

\slq{00:31:01}{In the case of geometric Langlands, we take $Y$ to be $\Loc(G)$ \dots\ birational to
$T^*\mathcal{M}(G)$. So $\mathcal{A}=\Hom(\Bcc,\Bcc)$ is the sheaf of holomorphic differential operators
on $\mathcal{M}(G)$. If $\mathcal{B}$ is any other brane, then $\Hom(\mathcal{B},\Bcc)$ is going to be a
module for $\mathcal{A}$ \dots\ In other words, the category on the ``automorphic'' side \dots\ is a
category of, roughly speaking, $\mathcal{D}$-modules on $\mathcal{M}(G)$.}

\subsection{为什么要四维版本；$\Bcc$ 的对偶是 opers}

\slq{00:32:23}{Mathematically, it is important to work with $\mathcal{D}$-modules on the stack $\Bun$,
not on a finite-dimensional moduli space $\mathcal{M}(G)$. Physically, what that means is that it is
important that a version of $\Bcc$ can be defined directly in four dimensions \dots\ (D.~Gaiotto and I
found an improved version \dots\ ``Knot Invariants From Four-Dimensional Gauge Theory,'' 2010.) That also
enabled us to answer a question that had been unclear and will be important later in this lecture:
what is the dual of $\Bcc$? \dots\ its dual is the structure sheaf of the variety of opers, a very
special Lagrangian submanifold $L_{\op}\subset\Loc(G^{\vee})$.}

\zm{这是全场的枢纽句。注意三件事：(i) 数学上必须用 stack $\Bun$，物理上的对应物是"四维里直接定义的 $\Bcc$"；
(ii) 对偶把 $A$-brane 送到 $B$-brane，所以 $\Bcc$ 的对偶是 $\Loc(G^\vee)$ 上的相干层；
(iii) 具体是哪一层，Beilinson–Drinfeld 早已给出：opers 的结构层。Witten 明说
"a gauge theory explanation of that fact was initially unclear, at least to me"。}

\subsection{小结（讲者自己的五条）}

\slq{00:34:21}{In summary the main ideas in the gauge theory/geometric Langlands correspondence are
\begin{itemize}\setlength\itemsep{0pt}
\item electric-magnetic duality $G\leftrightarrow G^{\vee}$ of supersymmetric gauge theory in four dimensions
\item ``twisting'' to make a dual pair of topological field theories
\item the dual theories have dual sets of line operators, with 't Hooft operators that are related to
geometric Hecke transformations, and their dual Wilson operators
\item on one side of the duality there is a distinguished brane $\Bcc$, establishing a map from
$A$-branes to $\mathcal{D}$-modules on $\Bun$
\item on the other side, one is studying $B$-branes on $\Loc(G^{\vee})$ in the complex structure in
which it parametrizes flat bundles
\end{itemize}}

\section{从形变量子化到量子化：Gukov–Witten}

\subsection{差别在哪：$S^2$ 的例子}

\slq{00:35:15}{This story so far involves deformation quantization, not quantization. Let me pause to
underline the difference, with the concrete example of a two-sphere $x^2+y^2+z^2=j^2$ viewed as a
symplectic manifold with its usual rotation-invariant symplectic form ($\omega=\mathrm{d}x\,\mathrm{d}y/z$).
In deformation quantization, we start with the commutative algebra of functions $\mathbb{C}[x,y,z]$ and
we want to deform it to a noncommutative algebra. \dots\ In the present case, we can just write down the
answer, which is given by the $\mathfrak{su}(2)$ Lie algebra
\[
[x,y]=-\mathrm{i}\hbar z,\qquad [y,z]=-\mathrm{i}\hbar x,\qquad [z,x]=-\mathrm{i}\hbar y .
\]}

\zm{这一段是整场报告的"概念分界线"。形变量子化只保证 $\{\ ,\ \}$ 的\textbf{一阶}对上；
量子化还要求有一个\textbf{Hilbert 空间}、有内积、有正定性、有真正的谱。
报告后半段的所有麻烦（$\tau$、$\Theta_\tau$、正定性、实 opers）都来自这个区别。}

\subsection{Gukov–Witten 判据}

\slq{00:40:44}{S.~Gukov and I (``Branes and Quantization,'' 2008) asked whether the $A$-model in this
general setup can describe quantization, and not just deformation quantization. The answer is yes, under
certain conditions. \dots\ we are discussing a complex symplectic manifold $Y$ \dots\ viewed as a real
symplectic manifold with symplectic form $\omega=\operatorname{Im}\Omega$. Let us discuss Lagrangian
$A$-branes, that is branes supported on a submanifold $L\subset Y$ that is Lagrangian for
$\operatorname{Im}\Omega$.}

\slq{00:42:34}{Gukov and I considered the opposite case of an $A$-brane $\mathcal{B}$ whose support $L$ is
Lagrangian for $\operatorname{Im}\Omega$ -- as it must be -- but is \emph{symplectic} for
$\operatorname{Re}\Omega$. We argued that in this case $\mathcal{H}=\Hom(\mathcal{B},\Bcc)$ represents a
quantization of $L$ with symplectic structure $\operatorname{Re}\Omega$. \dots\ $\mathcal{H}$ is an
$\mathcal{A}$-module which one can show has appropriate properties to be a quantization of $M$.}

\zm{判据可以写成一句：\textbf{量子化 = 挑一个"只对 $\operatorname{Im}\Omega$ 是 Lagrangian"的 brane $\mathcal{B}$}，
$\operatorname{Re}\Omega$ 在 $L$ 上仍是辛的，于是 $L$ 自己带了要量子化的辛结构，而 $\mathcal{H}=\Hom(\mathcal{B},\Bcc)$
自动是 $L$ 上的量子化。注意这里把"$L$ 的对偶是 $M$"用上了：$M$ 是 $L$ 的复化里的实形式。}

\subsection{什么时候是 Hilbert 空间：反全纯对合 $\tau$}

\slq{00:44:57}{I have actually omitted so far a key point: what structure is needed so that
$\mathcal{H}=\Hom(\mathcal{B},\Bcc)$ is a Hilbert space, not just a vector space? In general \dots\
there is a nondegenerate bilinear (not hermitian) pairing
\[
\Hom(\mathcal{B}_1,\mathcal{B}_2)\otimes\Hom(\mathcal{B}_2,\mathcal{B}_1)\to\mathbb{C}.
\]
To get a hermitian inner product on $\Hom(\mathcal{B}_1,\mathcal{B}_2)$ is therefore the same as finding
an antilinear map from $\Hom(\mathcal{B}_2,\mathcal{B}_1)$ to $\Hom(\mathcal{B}_1,\mathcal{B}_2)$.
Gukov and I showed that there is such an antilinear map, which I will call $\Theta_\tau$, if there is an
antiholomorphic map $\tau:Y\to Y$ with $M$ as a component of its fixed point set.}

\slq{00:48:54}{Given $\tau$, the hermitian product on $\Hom(\mathcal{B}_1,\mathcal{B}_2)$ for any pair of
$\tau$-invariant branes is defined by
\[
\langle\psi,\chi\rangle=(\Theta_\tau\psi,\chi)
\]
where $\Theta_\tau$ is the antilinear mapping that is defined using $\tau$, and $(\ ,\ )$ is the bilinear
pairing of the $A$-model.}

\zm{为什么 $\tau$ 是对的条件：若全纯函数在 $\mathcal{H}$ 上作用，而 $\mathcal{H}$ 是 Hilbert 空间，
就必须能判定"哪些全纯函数是自伴的"。$\tau$ 把 $Y$ 的实点集 $M$ 固定住，于是"在 $M$ 上取实值"
给出一个自然的自伴判定。这也是后面"实 opers"一词的来历。}

\subsection{与几何量子化的比较}

\slq{00:51:01}{A limitation of my paper with Gukov is that we did not understand very much about how to
compare brane quantization to geometric quantization. \dots\ if $M$ has a polarization $\mathcal{P}$ in
the sense of geometric quantization, and a complexification $Y$ that is suitable for quantization by
branes, and if $\mathcal{P}$ analytically continues to what we call a holomorphic polarization of $Y$,
then the two methods to quantize $M$ agree.}

\zm{这是报告里唯一一处"两套量子化理论等价"的明确陈述，条件很强（极化要能解析延拓）。
EFK 的 Hilbert 空间正是几何量子化那一侧，所以这句话是整场报告的技术前提：\textbf{它允许我们把
EFK 的 $\mathrm{L}^2$ 换成 brane quantization 的 $\Hom(\mathcal{B},\Bcc)$，而不损失任何东西。}}

\section{EFK 的 Hilbert 空间与 Hitchin 谱}

\subsection{两边的对象}

\slq{00:52:00}{They considered a Hilbert space $\mathcal{H}$ of $\mathrm{L}^2$ functions (or better,
half-densities) on $\mathcal{M}(G)$ \dots\ They constructed operators on $\mathcal{H}$ that are related
to the usual constructions of geometric Langlands, and found interesting duality theorems and
conjectures concerning the action of these operators.}

\slq{00:52:37}{In geometric quantization, one would understand the $\mathrm{L}^2$ functions on
$\mathcal{M}(G)$ in terms of quantization of the cotangent bundle $T^*\mathcal{M}(G)\cong\Loc(G)$.
In other words, the Hilbert space of EFK is what one gets if [one] takes the Higgs bundle moduli space
$\Loc(G)$ with real symplectic structure $\omega=\operatorname{Im}\Omega$, and quantizes it via geometric
quantization, using the fact that it is a cotangent bundle.}

\zm{所以"EFK 的 $\mathcal{H}$"与"$\Loc(G)$ 的 brane 量子化"是同一个东西的两种描述。
Witten 后面不再区分，直接写 $\mathcal{H}=\Hom(\mathcal{B},\widehat{\Bcc})$。}

\subsection{可观测量：$\mathcal{A}\otimes\overline{\mathcal{A}}$}

\slq{00:57:07}{The Higgs bundle moduli space has a Hitchin fibration $\pi:\Loc\to B$, where
$B\simeq\mathbb{C}^n$ for some $n$. The linear functions on $B$ are Hitchin's commuting Hamiltonians.
Classically the global holomorphic functions on $\Loc$ are just the pullbacks of functions on $B$,
i.e.\ the algebra $\mathcal{A}_0$ of holomorphic functions on $\Loc$ is the algebra of polynomial
functions of the Hitchin Hamiltonians.}

\slq{00:57:46}{For brane quantization, we define the canonical coisotropic brane $\widehat{\Bcc}$ of
$\widehat{Y}=Y_1\times Y_2$. It is just a product $\widehat{\Bcc}=\mathcal{B}_{\mathrm{cc},1}\times\mathcal{B}_{\mathrm{cc},2}$ of canonical
coisotropic branes on the two factors. So the algebra that acts on the quantization of $Y$ is
$\Hom(\widehat{\Bcc},\widehat{\Bcc})=\mathcal{A}\otimes\overline{\mathcal{A}}$.}

\zm{$Y_1$ 与 $Y_2$ 是 $\Loc$ 的两种复结构（同一个实流形，全纯与反全纯），$\mathcal{A}$ 与 $\overline{\mathcal{A}}$
是形变量子化后的全纯/反全纯函数代数。$\mathcal{A}\otimes\overline{\mathcal{A}}$ 里两类算子\textbf{平凡交换}——
真正的信息在"它们的联合谱是什么"，而这只能靠对偶来算。}

\subsection{对偶怎么作用：三个 brane 的对偶}

\slq{01:01:23}{To apply duality to $\mathcal{H}=\Hom(\mathcal{B},\widehat{\Bcc})$ where
$\widehat{\Bcc}=\mathcal{B}_{\mathrm{cc},1}\times\mathcal{B}_{\mathrm{cc},2}$, naively we need to understand the duals of the three branes involved
\dots\ $\mathcal{B}_{\mathrm{cc},1}$ is the brane associated to deformation quantization of the ring of holomorphic functions,
and \dots\ its dual is the structure sheaf of the variety of opers, $L_{\op}\subset\Loc(G^{\vee})$.
\dots\ $\mathcal{B}_{\mathrm{cc},2}$ is the brane associated to deformation quantization of the ring of antiholomorphic
functions. So its dual is the structure sheaf of $L_{\opb}$, the Lagrangian submanifold that parametrizes
flat bundles whose antiholomorphic structure obeys the oper condition.}

\zm{注意 $L_{\op}$ 与 $L_{\opb}$ 都在 $\Loc(G^\vee)$ 的\textbf{同一个}复结构里（用 $G^\vee_{\mathbb{C}}$ 的全纯结构），
所以才能谈"交集"。讲者 01:07:03 专门解释了这一点：

\slq{01:07:03}{A flat bundle has both a holomorphic structure and an antiholomorphic structure, both of
which vary holomorphically. So in fact there is a holomorphic map
$\Loc(G,C)\to\Bun(C)\times\Bun(\overline{C})$, where $\overline{C}$ is $C$ with opposite complex
structure. A flat bundle $E$ (for $SL_2$, say) is an oper holomorphically if holomorphically it is a
nontrivial extension
\[
0\to K^{1/2}\to E\to K^{-1/2}\to 0
\]
and it is an oper antiholomorphically if antiholomorphically it is a nontrivial extension
\[
0\to\overline{K}^{1/2}\to E\to\overline{K}^{-1/2}\to 0 .
\]}
}

\subsection{对角 brane 与 unfolding trick}

$\mathcal{B}$ 支撑在 $Y_1\times Y_2$ 的对角线上（"$\Loc$ 的两种复结构对齐"）。看起来很麻烦，
但报告给出一个技巧绕开它：

\slq{01:03}{There is an ``unfolding trick'' that shows we do not have to worry about dualizing the
diagonal: \dots\ So we just dualize $\Bcc$ and $\widetilde{\Bcc}$ to get $L_{\op}$ and $L_{\opb}$.}

\zm{unfolding 的图（01:05 的两张）把"沿对角线粘起来的带子"展平成"一整块矩形"，
于是对角线 brane 的对偶就把 $\Hom$ 变成 $\Hom(L_{\opb},L_{\op})$。这张图是整场报告最"物理"的一步，
也是代数几何那边不一定能直接看到的步骤。}

\subsection{正定性 $\Leftrightarrow$ 实 opers}

\slq{01:09:54}{By the method I explained before using an antiholomorphic involution $\tau$,
$\Hom(\widetilde{\Bcc},\Bcc)$ carries a hermitian inner product, which we expect to be positive as it
corresponds to the quantization of a cotangent bundle. The dual $\Hom(L_{\opb},L_{\op})$ likewise
carries a nondegenerate inner product \dots\ However, for generic $B$-branes this inner product is not
positive-definite. Conditions that make the hermitian form positive-definite were formulated by EFK,
who showed that these conditions are satisfied for $GL_1$ (by a direct but surprisingly non-trivial
computation) and for $SL_2$ (by a theorem of Faltings).}

\slq{01:11:20}{Consider a point $x\in L_{\op}\cap L_{\opb}$, corresponding to a flat bundle $E$ that is an
oper both holomorphically and antiholomorphically. The antiholomorphic involution $\tau^{\vee}$ \dots\
maps $E$ to the complex conjugate flat bundle $\overline{E}$ \dots\ If $E$ is not isomorphic to
$\overline{E}$, then they are both null vectors \dots\ But the duality predicts that it should be
positive-definite, so we expect that, as conjectured by EFK and proved in some cases,
$E$ is always isomorphic to $\overline{E}$. Thus \dots\ the flat $G_{\mathbb{C}}$ bundles that are opers
both holomorphically and antiholomorphically are actually real.}

\zm{这一段是全报告最漂亮的"物理逼出数学"：正定性 $\Rightarrow$ 交点必须成对出现且不能是零向量
$\Rightarrow$ $E\cong\overline{E}$ $\Rightarrow$ 结构群可降到实形式。EFK 的猜想在这里变成一个逻辑必然，
条件是"对偶成立 + 内积正定"。}

\slq{01:12:45}{Assuming the intersection points are all real in that sense, $\Hom(L_{\opb},L_{\op})$
is positive-definite if and only if all intersection points of $L_{\opb}$ and $L_{\op}$ are always
transverse. \dots\ Bundles with this property are what EFK call \emph{real opers}.}

\subsection{联合谱}

\slq{01:14:41}{Thus assuming the duality is true and the hermitian form is positive-definite, the joint
spectrum of $\mathcal{A}$ and $\overline{\mathcal{A}}$ corresponds to real opers. \dots\ according to
Beilinson and Drinfeld, an element $x$ of $\mathcal{A}$ (1) is a holomorphic differential operator on
$\Bun$ and (2) corresponds to a holomorphic function $f_x$ on $L_{\op}$. Similarly \dots\ an element
$x'$ of $\overline{\mathcal{A}}$ \dots\ corresponds to a holomorphic function $f_{x'}$ on $L_{\opb}$.
So an intersection point $p$ of $L_{\op}$ and $L_{\opb}$ determines a pair of eigenvalues of the
differential operators $x,x'$, namely $f_x(p)$ and $f_{x'}(p)$. This is the proposal for the joint
spectrum of Hitchin's quantized Hamiltonians.}

\zm{这就是 EFK 那篇的"谱"定理在物理语言里的样子：\textbf{联合谱 = 实 opers 上的函数值}。
数学上它等价于"$\mathcal{A}$ 是 $\Bun$ 上的全纯微分算子代数，且其在 $L_{\op}$ 上的特征值由 $p$ 给出"；
物理上它只是"对偶把 $\mathcal{A}$ 变成 $L_{\op}$ 上的函数"。}

\section{Hecke 算子的量子化与对偶预测}

\subsection{交换性来自"线可以滑动"}

\slq{01:22:43}{A similar argument shows that Hecke operators $T_R(p)$ commute with the quantized version
of Hitchin's holomorphic (or antiholomorphic) Hamiltonians.}

\zm{"similar argument" 指 01:21 那几页：两个线算子在不同点上可以互相滑过去（拓扑不变），
所以对易。和 00:13 那句 "they commute, but a little less obviously"（同一点上不同表示）对照着读，
能看出 Witten 反复用同一张图论证三件不同的事：$T$ 与 $T$ 交换、$T$ 与 $W$ 交换、$T$ 与 $H$ 交换。}

\subsection{Wilson 算子在 $G^\vee$ 侧怎么变成函数}

\slq{01:27:32}{Given a $G^{\vee}$ bundle $E$ with connection over $\Sigma\times C$, and a representation
$R$ of $G^{\vee}$, we form the associated bundle $E_R=E\times_{G^{\vee}_{\mathbb{C}}}R$, which also comes
with a connection. $W_R(\gamma)$, for a path $\gamma$ in $\Sigma\times C$, is defined by parallel transport
of the connection on $E_R$ along $\gamma$. For the present application, we fix a point $p\in C$ and two
points $a,b$ on the right and left boundaries of $\Sigma$, and a path in $\Sigma\times p$ from $a\times p$
to $b\times p$ \dots\ parallel transport defines, for each connection, a linear transformation
$W_R(p):E_{R,a\times p}\to E_{R,b\times p}$, or equivalently an element
\[
W_{R,p}\in\Hom(E_{R,a\times p}\otimes E_{R',b\times p},\mathbb{C}),
\]
where $R'$ is the dual representation to $R$.}

\slq{01:29:48}{A quantum operator is going to come from something that is a complex valued function of
connections. We do not have that yet \dots\ To get an operator, we need to supply elements
$v\in E_{R,a\times p}$, $w\in E_{R',b\times p}$ and then $W_{R,p}(v\otimes w)$ is the function of
connections whose quantization will be an operator.}

\zm{这里的逻辑缺口是"$W_{R,p}$ 依赖边界上的纤维，而纤维不是数值"。补法是\textbf{用边界的几何挑出纤维里的向量}，
即下面的 "corners"。}

\subsection{Corners 与自旋结构无关性}

\slq{01:30:44}{the boundary condition on the right boundary is that the $G^{\vee}$ bundle $E$, restricted
to the right boundary, is an antiholomorphic oper. This means that there is a nonsplit exact sequence
\[
0\to\overline{K}^{1/2}\xrightarrow{\overline{j}}E_R\to\overline{K}^{-1/2}\to 0
\]
\dots\ So if we pick a vector $v\in\overline{K}_p^{1/2}$ (the fiber of $\overline{K}^{1/2}$ at $p$) then
this gives a vector $\overline{j}(v)\in E_{R,a\times p}$.}

\slq{01:31:28}{Similarly on the left boundary the boundary condition says that $E$ is a holomorphic oper,
implying that there is a nonsplit exact sequence
\[
0\to K^{1/2}\xrightarrow{j}E_{R'}\to K^{-1/2}.
\]
So we pick $w\in K_p^{1/2}$ \dots\ In other words, the ``corners'' correspond to
$v\in\overline{K}^{1/2}$ and $w\in K^{1/2}$. Once they are picked, we define
\[
\widehat{W}_{R,p}=W_{R,p}\bigl(\overline{j}(v)\otimes j(w)\bigr)
\]
and this is the complex-valued function of connections that we will interpret as a quantum operator.}

\slq{01:32:20}{This step is trivial, because in the $B$-model \dots\ one can actually assume that the
connection is pulled back from $C$. $\widehat{W}_{R,p}$ is just the natural dual pairing
$(\ ,\ ):E_{R,p}\otimes E_{R',p}\to\mathbb{C}$, and therefore its value at a given real oper is simply
\[
(\overline{j}(v),j(w)).
\]
\dots\ it actually only depends on $v\otimes w\in K^{1/2}\otimes\overline{K}^{1/2}=|K|$, and thus in
particular as observed by EFK there is no dependence on a choice of spin structure.}

\subsection{对偶给出的预测}

\slq{01:33:46}{The duality then predicts that also on the dual side, the definition of $\mathcal{T}(p)$
as an operator requires ``corners'' $v,w\in\overline{K}_p^{1/2},K_p^{1/2}$. This is true, as shown by
EFK using an algebro-geometric formula. The duality further predicts that the eigenvalues of the
resulting operator are what we found on the $G^{\vee}$ side: they are
\[
(\overline{j}(v),j(w))
\]
for all the possible real opers. (As a result of a lecture by E.~Frenkel two hours ago, [I] learned of a
slight refinement. Because $SL_2$ has a nontrivial center $\{\pm1\}$, parallel transport from left to
right is uniquely determined only up to sign. Hence actually the eigenvalues are
\[
\pm(\overline{j}(v),j(w))
\]
with both signs occurring.)}

\zm{这句"两小时前听了 Frenkel 的报告才知道要加 $\pm$"是现场感最强的一处：说明这条预测在当时仍在打磨。
$\pm$ 的来源是中心 $Z(SL_2)=\{\pm1\}$ 使平行移动只在中心元意义下确定。}

\section{讲者自己留下的问题}

按报告里明确的措辞（"I don't know" / "we don't understand" / "unclear"）整理：

\begin{enumerate}[label=\arabic*.]
\item \textbf{正定性只在两个特例已知}（01:09:54）：$GL_1$ 靠直接计算，$SL_2$ 靠 Faltings 的定理；
一般 $G$ 没有一般性论证。
\item \textbf{实 opers 是否总落在分裂实形式上}（01:11:20）：Witten 原话 "I don't know if it is always
the split real form."
\item \textbf{交点横截性}（01:12:45）：把正定性化归为"$L_{\op}$ 与 $L_{\opb}$ 的交点总是横截"，
但这本身是 EFK 的结果/猜想，报告未证。
\item \textbf{与几何量子化的比较只在极化可解析延拓时成立}（00:51:01）：这是全报告的技术前提，不是定理。
\item \textbf{$\Bcc$ 的对偶为什么是 opers}：数学上是 Beilinson–Drinfeld，物理上 "initially unclear, at least to me"（00:32:23）。
\item \textbf{一般 Lagrangian brane 的对偶}：只有当 $L$ 来自 hyperholomorphic 的点（如 Hitchin 纤维）时才是复 Lagrangian；
一般情况下 $L$ 只对 $\operatorname{Im}\Omega$ Lagrangian，那里的量子化行为报告未展开（00:40:44）。
\end{enumerate}

\section{术语与记号}

\begin{center}\small
\begin{tabular}{@{}p{4.6cm}p{10.6cm}@{}}
\toprule
记号 / 术语 & 含义（本笔记用法）\\
\midrule
$\mathcal{N}=4$ SYM & 四维 $\mathcal{N}=4$ 超对称 Yang–Mills，报告的全部物理起点\\
$\tau,\ \tau^{\vee}$ & 复耦合 $\frac{\theta}{2\pi}+\frac{4\pi}{e^2}$ 与对偶群侧的对应量，$\tau^\vee=-1/(n_{\mathfrak{g}}\tau)$\\
$\mathcal{C},\ \mathcal{C}^{\vee}$ & 固定 $C$ 后，$G$ 侧与 $G^{\vee}$ 侧的边界条件范畴\\
$T(p),\ T_R(p)$ & 几何 Hecke 函子/算子（$G$ 侧，来自 't Hooft 线算子）\\
$W(p),\ W_R(p)$ & Wilson 函子/算子（$G^{\vee}$ 侧），满足 $W_RW_{R'}=W_{R\otimes R'}$\\
$\Loc(G)$ & $G$-Higgs 丛模空间，$\simeq T^*\mathcal{M}(G)$（双有理）\\
$\Bcc$ & canonical coisotropic $A$-brane，支撑全空间，$\Hom(\Bcc,\Bcc)\sim$ 形变量子化\\
$\widehat{\Bcc}$ & $Y_1\times Y_2$ 上的 $\mathcal{B}_{\mathrm{cc},1}\times\mathcal{B}_{\mathrm{cc},2}$，其自同态为 $\mathcal{A}\otimes\overline{\mathcal{A}}$\\
$L_{\op},\ L_{\opb}$ & opers / 反全纯 opers 构成的（复）Lagrangian 子流形，位于 $\Loc(G^{\vee})$\\
real oper & 既是全纯又是反全纯 oper 的平坦丛，且结构群可降到实形式\\
$\tau$（对合） & $Y$ 的反全纯对合，不动点集含 $M$；用来造 $\Theta_\tau$\\
$\Theta_\tau$ & 由 $\tau$ 定义的反线性映射，把双线性配对变成 Hermite 内积\\
corners & 边界上用来挑出 $E_R$ 纤维里向量的数据 $v\in\overline{K}^{1/2}_p,\ w\in K^{1/2}_p$\\
\bottomrule
\end{tabular}
\end{center}

\section{时间轴}

\begin{center}\small
\begin{tabular}{@{}llp{9.6cm}@{}}
\toprule
时间 & 块 & 内容\\
\midrule
00:00:30 & 3 & 与 Gaiotto 合作，目标是用 QFT 理解 EFK（arXiv:1908.09677）\\
00:01:04 & 3 & $\tau$ 与电-磁对偶 $\tau^{\vee}=-1/(n_{\mathfrak{g}}\tau)$\\
00:04:50 & 6 & 场论给 $d$-流形配数、给 $(d-1)$-流形配范畴\\
00:06:49 & 7 & $\Hom(\mathcal{B}_1,\mathcal{B}_2)$ 的定义；自同态是代数\\
00:10:00 & 9–12 & 线算子 = 边界条件范畴的自函子；Wilson vs 't Hooft\\
00:15:35 & 14 & $T(p)/W(p)$ 与几何 Satake：$W_RW_{R'}=W_{R\otimes R'}$\\
00:20:00 & 15 & 几何 Langlands = $\Loc(G)$ 的 $A$-模型与 $\Loc(G^{\vee})$ 的 $B$-模型的镜像对称\\
00:29:01 & 17 & $\Bcc$ 的定义；$Y=T^*W$ 时 $\Hom(\Bcc,\Bcc)$ 是全纯微分算子\\
00:31:01 & 19 & 取 $Y=\Loc(G)$，得到 $\Bun$ 上的 $\mathcal{D}$-模\\
00:32:23 & 20 & $\Bcc$ 的对偶 = opers 的结构层（Beilinson–Drinfeld）\\
00:34:21 & 21 & 小结五条\\
00:35:15 & 21 & $S^2$ 例子：形变量子化 vs 量子化\\
00:40:44 & 29 & Gukov–Witten 判据：$L$ 对 $\operatorname{Re}\Omega$ 辛 $\Rightarrow$ $\Hom(\mathcal{B},\Bcc)$ 是量子化\\
00:44:57 & 31 & 需要 Hilbert 结构；反全纯对合 $\tau$ 与 $\Theta_\tau$\\
00:48:54 & 35 & $\langle\psi,\chi\rangle=(\Theta_\tau\psi,\chi)$\\
00:51:01 & 39 & 与几何量子化一致的条件（极化解析延拓）\\
00:52:37 & 41 & EFK 的 $\mathrm{L}^2(\mathcal{M}(G))$ = 量子化 $\Loc(G)=T^*\mathcal{M}(G)$\\
00:57:07 & 47 & Hitchin 纤维化与哈密顿量；可观测量 $\mathcal{A}\otimes\overline{\mathcal{A}}$\\
01:01:23 & 52 & 三个 brane 的对偶：$\mathcal{B}_{\mathrm{cc},1}\mapsto L_{\op}$，$\mathcal{B}_{\mathrm{cc},2}\mapsto L_{\opb}$\\
01:05:22 & 54 & unfolding trick，只对偶 $\Bcc$ 与 $\widetilde{\Bcc}$\\
01:07:03 & 58 & 为什么两者在同一复结构里；oper 的扩张 $0\to K^{1/2}\to E\to K^{-1/2}\to 0$\\
01:09:54 & 59 & 正定性；$GL_1$ 与 $SL_2$（Faltings）的特例\\
01:11:20 & 60 & 正定性 $\Rightarrow$ $E\cong\overline{E}$ $\Rightarrow$ real opers\\
01:14:41 & 62 & 联合谱 = real opers 上 $f_x(p),f_{x'}(p)$\\
01:22:43 & 69 & Hecke 算子与量子化 Hitchin 哈密顿量交换\\
01:27:32 & 75 & $W_R(p)$ 作为平行移动；$W_{R,p}\in\Hom(E_{R,a}\otimes E_{R',b},\mathbb{C})$\\
01:30:44 & 77 & 右边界 = 反全纯 oper，$v\in\overline{K}^{1/2}_p$\\
01:31:28 & 78 & 左边界 = 全纯 oper，$w\in K^{1/2}_p$；$\widehat{W}_{R,p}=W_{R,p}(\overline{j}(v)\otimes j(w))$\\
01:32:20 & 79 & 只依赖 $v\otimes w\in|K|$，无自旋结构依赖；值为 $(\overline{j}(v),j(w))$\\
01:33:46 & 80 & 对偶预测 Hecke 特征值；$SL_2$ 加 $\pm$（Frenkel 两小时前的报告）\\
\bottomrule
\end{tabular}
\end{center}

\section{方法与可信度（请先读这一段再引用）}

\textbf{这份笔记是怎么来的。}源视频无字幕、无自动字幕；讲义页直接从画面转录（逐张读图、公式转 LaTeX），
音轨用 faster-whisper small.en 分块转写并用术语表纠正。两层按时间轴合并，页与块的边界只依赖时间，
不依赖页码（重跑时页码会在 201$\leftrightarrow$208 间漂移）。

\textbf{哪些地方可能错。}按风险从高到低：
\begin{enumerate}[label=\arabic*.]
\item \textbf{公式的上下标}：转录时最易错的形态是"看起来对但差一个指标"。本笔记抽核了 18 处
（$\tau^{\vee}$、$n_{\mathfrak{g}}$、$\Theta_\tau$、$\overline{K}^{1/2}$ 的指数、oper 扩张的方向、
$W_{R,p}$ 的 $\Hom$ 参数），未发现错误，但\textbf{没有做系统抽核}，所以定为 B 级。
\item \textbf{OCR 层面的小词}：例如 "if takes the Higgs bundle" （原文如此，疑为 "if one takes"），
笔记中保留原样，不作静默修正。
\item \textbf{口述转写的同音词}：已用术语表纠正（branes/brains、oper/au pair、Hitchin/Hitchen 等），
但人名与冷僻词仍可能有误。
\item \textbf{讲义跳页}：228 页中讲者只翻到 201 页，缺的 27 页没有任何信息，本笔记无法补。
\end{enumerate}

\textbf{怎么自己复核。}最省事的办法是看随附的 \texttt{source/*.jsonl}：每条都有 \texttt{t\_start} 与
\texttt{frame}，回到视频对应时刻即可对照。发现问题请走仓库 issue，并在该份 \texttt{README.md} 的
"复核"行里记录。

\vfill
\noindent\rule{\textwidth}{0.4pt}\\[2pt]
{\footnotesize 本笔记遵循《视频讲义提取与笔记规范 v1.0》（\texttt{04\_playbook/24\_视频讲义提取与笔记规范\_v1.md}）。
讲义原文与公式的内容权利归讲者 Edward Witten 及会议主办方；中文叙述、结构、时间轴为本仓自产，可按 CC BY 4.0 使用。
如权利人要求，本目录整体撤下。}

\end{document}
